% Part: applied-modal-logics
% Chapter: epistemic-logic
% Section: properties-accessibility

\documentclass[../../../include/open-logic-section]{subfiles}

\begin{document}

\olfileid{aml}{el}{acc}

\olsection{علاقات الإتاحة وأصول المعرفة}

قد عُرفت مطابقة الأطر في المنطقيات الجهوية المعتادة؛ فلننظر الآن في معنى
ال!!{formula}s المميزة إذا أولت تأويلًا معرفيًا. وقد تقدم أن منطق المعرفة
يفسر عادةً في S5؛ والمطلوب بيان كيف تمثل أصول المعرفة المختلفة، وما يقابل
كل أصل منها من شروط الإطار.

فقد كانت ال!!{formula}s الجهوية، في المنطق الجهوي المعتاد، تميز خواص
متباينة لعلاقات الإتاحة. وفي الجدول الآتي طائفة منها، بإزاء ما يقابلها من
أصول المعرفة.

\begin{table}[t]
    \begin{tabular}{| p{.48\textwidth} || p{.48\textwidth} |}
      \hline
      {\emph{إذا كانت $R$ \dots}} & {\emph{فإن \dots صادقة في~$\mModel{M}$:}} \\
      \hline \hline
        
      & $\Knows (p \lif q) \lif (\Knows p \lif \Knows q)$ 
      \hfill \newline{(الانغلاق)} \\
      \hline
      \emph{انعكاسية}: $\forall w Rww$  
      & $\Knows p \lif p$ \hfill \newline{(الصدقية)} \\
      \hline
      \emph{متعدية}: \newline
      $\forall u \forall v \forall w ((Ruv \land Rvw) \lif Ruw)$ & 
      $\Knows p \lif \Knows \Knows p$ \hfill 
           \newline{(الاستبطان الإيجابي)} \\
      \hline 
      \emph{إقليدية}: \newline
      $\forall w \forall u \forall v ((Rwu \land Rwv) \lif Ruv)$ &  
      $\lnot \Knows p \lif \Knows \neg \Knows p$ \hfill 
        \newline{(الاستبطان السلبي)} \\
      \hline
    \end{tabular}
    \caption{أربعة مبادئ معرفية.}
    \ollabel{tab:four}
  \end{table} 

أما الصدقية، المقابلة لمسلمة~$T$، فهي أقل هذه الأصول خلافًا؛ إذ تقضي بأن
!!a{formula} المعلومة لا بد أن تكون صادقة. وأما الانغلاق والاستبطان
الإيجابي والاستبطان السلبي ففيها خلاف أشد.

والانغلاق، المقابل لمسلمة~$K$، معناه أن معرفة الفاعل مغلقة تحت الاستلزام.
وقد يقرب هذا في بعض الأمثلة: فأنا أعلم أن وجودي في فيكتوريا يستلزم وجودي
في جزيرة فانكوفر، وأعلم --- إذا صرفنا النظر عن فروض الشك الغريبة --- أنني
في فيكتوريا؛ فيظهر أنني أعلم أنني في جزيرة فانكوفر. غير أنا لا نستخرج
دائمًا جميع لوازم معلوماتنا، فيفضي الأصل في أحوال أخرى إلى نتائج أقل
بداهة.

ويسمى الاستبطان الإيجابي أحيانًا أصل~KK: إذا علمت شيئًا علمت أنني أعلمه؛
وهو النظير المعرفي لمسلمة~4. وعلى قياسه يقول الاستبطان السلبي: إذا كنت
\emph{لا أعلم} شيئًا علمت أنني لا أعلمه؛ وهو نظير مسلمة~5. ولكل واحد
منهما فيما يظهر أمثلة مضادة مألوفة: فقد أعلم شيئًا في الواقع وأنا غير
متيقن أنني أعلمه.


\end{document}
